\documentclass[11pt]{article}
\usepackage[margin=1in]{geometry}
\usepackage{listings}
\usepackage{url}
\usepackage[colorlinks=true,linkcolor=blue,citecolor=blue,urlcolor=blue]{hyperref}
\lstset{basicstyle=\small\ttfamily, columns=fullflexible, breaklines=true,
        frame=single, framerule=0.4pt, backgroundcolor=\color[gray]{0.97}}

\title{Modula-3 / C++ Interoperability:\\Design Notes from a Toy Binding}
\author{Mika Nystr\"om}
\date{May 8, 2026}

\begin{document}
\maketitle
\tableofcontents
\clearpage

%% -------------------------------------------------------------------
\section{Introduction}

Modula-3 (M3) can call C functions via \texttt{EXTERNAL} declarations,
but C++ libraries have no C calling convention.  This document describes
a three-layer binding architecture demonstrated with a toy C++ class
hierarchy, and identifies the design decisions that generalize to
real libraries like Open3D, Eigen, or any C++ API.

Source code:
\begin{itemize}
\item \url{https://github.com/mikanystrom/intel-async/tree/al-ga-surface-analysis/async-toolkit/m3utils/m3utils/cxxbridge}
\end{itemize}

%% -------------------------------------------------------------------
\section{Architecture}

The binding consists of three layers:

\begin{center}
\begin{tabular}{lll}
\textbf{Layer} & \textbf{Language} & \textbf{Role} \\
\hline
C++ library  & C++  & The real implementation (\texttt{shapes.hpp}) \\
C bridge     & C++ (\texttt{extern "C"}) & Flat C API over C++ objects (\texttt{shapes\_c.h}) \\
M3 binding   & Modula-3 & Idiomatic M3 types and GC integration \\
\end{tabular}
\end{center}

The C bridge is compiled as a static library (\texttt{libtoyshapes.a})
and linked into the M3 program along with \texttt{-lc++} for the C++
runtime.

%% -------------------------------------------------------------------
\section{Object Lifetime}

\subsection{The problem}

C++ uses RAII: objects are destroyed when they go out of scope or when
\texttt{delete} is called.  M3 uses garbage collection: objects are
destroyed at an unpredictable time after they become unreachable.
When M3 holds a handle to a C++ object, who calls \texttt{delete}?

\subsection{Solution: weak-reference cleaners}

Each M3 wrapper object registers a \textbf{weak-reference cleaner}
at construction time via \texttt{WeakRef.FromRef}.  When the M3 GC
collects the wrapper, the cleaner calls the C destructor:

\begin{lstlisting}
REVEAL Shape = BRANDED REF RECORD
  handle : ADDRESS;  (* C++ object pointer *)
  owned  : BOOLEAN;  (* TRUE = we should free on GC *)
END;

PROCEDURE ShapeCleaner(READONLY w: WeakRef.T; r: REFANY) =
  VAR s := NARROW(r, Shape);
  BEGIN
    IF s.owned AND s.handle # NIL THEN
      ToyShapesC.ShapeFree(s.handle);
      s.handle := NIL;
    END;
  END ShapeCleaner;
\end{lstlisting}

\subsection{Ownership transfer}

When a Shape is added to a ShapeList, the C++ side takes ownership
of the handle.  The M3 side sets \texttt{owned := FALSE} to
suppress the destructor:

\begin{lstlisting}
PROCEDURE Add(sl: ShapeList; s: Shape) =
  BEGIN
    ToyShapesC.ShapeListAdd(sl.handle, s.handle);
    s.owned := FALSE;  (* C++ owns it now *)
    (* Keep M3 reference so GC doesn't collect the wrapper *)
    sl.members[sl.nMembers] := s;
    INC(sl.nMembers);
  END Add;
\end{lstlisting}

This is the M3 equivalent of \texttt{std::unique\_ptr::release()}.

%% -------------------------------------------------------------------
\section{Type Fidelity}

The C++ class hierarchy (\texttt{Shape}, \texttt{Circle},
\texttt{Rectangle}) is represented on the M3 side as a
\textbf{single opaque type} \texttt{Shape}.  Virtual dispatch
happens entirely on the C++ side---the M3 code calls
\texttt{ToyShapesC.ShapeArea(s.handle)}, which the C bridge
casts to \texttt{Shape*} and calls the virtual method.

An alternative would be to mirror the hierarchy in M3:
\begin{lstlisting}
TYPE Shape <: REFANY;
     Circle <: Shape;
     Rectangle <: Shape;
\end{lstlisting}
This would give M3 code compile-time type safety (e.g., calling
\texttt{Radius()} on a non-Circle would be a type error).  The
cost is maintaining a parallel type hierarchy and doing a
\texttt{TYPECASE} or \texttt{NARROW} at the bridge layer.

For this prototype, the single-type approach was chosen for
simplicity.  Real bindings should consider which approach better
serves the users.

%% -------------------------------------------------------------------
\section{String Marshaling}

C++ \texttt{std::string} cannot cross the C boundary directly.
The C bridge returns a \texttt{malloc}'d \texttt{char*}:

\begin{lstlisting}[language=C]
char *toy_shape_name(toy_shape_t s) {
    std::string n = static_cast<Shape*>(s)->name();
    char *r = (char*)malloc(n.size() + 1);
    if (r) strcpy(r, n.c_str());
    return r;
}
\end{lstlisting}

The M3 side converts to \texttt{TEXT} and frees the C buffer:

\begin{lstlisting}
PROCEDURE Name(s: Shape): TEXT =
  VAR cstr := ToyShapesC.ShapeName(s.handle);
  BEGIN
    t := M3toC.CopyStoT(cstr);
    Cstdlib.free(cstr);
    RETURN t;
  END Name;
\end{lstlisting}

\texttt{M3toC.CopyStoT} copies the bytes into a GC-managed
\texttt{TEXT}, so the \texttt{free()} is safe immediately after.

%% -------------------------------------------------------------------
\section{Callbacks}

\subsection{The problem}

C++ \texttt{std::function} captures state (closures), but C function
pointers do not.  The standard C pattern is a function pointer plus a
\texttt{void*} context:

\begin{lstlisting}[language=C]
typedef void (*toy_visitor_fn)(toy_shape_t s, size_t i, void *ctx);
void toy_shapelist_foreach(toy_shapelist_t sl,
                           toy_visitor_fn fn, void *ctx);
\end{lstlisting}

On the M3 side, we want to pass an object with a method (a closure):

\begin{lstlisting}
TYPE VisitorClosure = OBJECT METHODS
  visit(s: Shape; index: CARDINAL);
END;
\end{lstlisting}

\subsection{Solution: trampoline + LOOPHOLE}

The M3 side passes a \textbf{trampoline} procedure as the function
pointer, and the closure object's address as the context:

\begin{lstlisting}
PROCEDURE ForEach(sl: ShapeList; v: VisitorClosure) =
  BEGIN
    ToyShapesC.ShapeListForEach(
      sl.handle, VisitorTrampoline, LOOPHOLE(v, ADDRESS));
  END ForEach;

PROCEDURE VisitorTrampoline(h: Handle; index: CARDINAL;
                            ctx: ADDRESS) =
  VAR v := LOOPHOLE(ctx, VisitorClosure);
      s := NEW(Shape, handle := h, owned := FALSE);
  BEGIN
    v.visit(s, index);
  END VisitorTrampoline;
\end{lstlisting}

The trampoline has a fixed address (it is a top-level procedure,
not a closure), so it is safe to pass as a C function pointer.
The \texttt{LOOPHOLE} converts between the M3 traced reference
and the C \texttt{void*}.

\textbf{GC safety}: During the callback, the closure \texttt{v} must
not be collected or moved.  Since \texttt{ForEach} holds a reference
to \texttt{v} on the M3 stack, the GC will not collect it.  CM3's
conservative stack scanning ensures that the traced reference is
seen even though we passed it as an untraced \texttt{ADDRESS}.

\subsection{Closures with state}

The predicate closure demonstrates captured state---the point
coordinates \texttt{px}, \texttt{py} are fields of the closure
object:

\begin{lstlisting}
TYPE ContainsPred = PredicateClosure OBJECT
  px, py: LONGREAL;
OVERRIDES
  test := ContainsTest;
END;

(* Create and use: *)
VAR pred := NEW(ContainsPred, px := 1.0d0, py := 0.0d0);
    hits := ToyShapes.Filter(sl, pred);
\end{lstlisting}

The state travels through the C bridge via the \texttt{void*}
context without the C code knowing anything about it.

%% -------------------------------------------------------------------
\section{Exception Handling}

C++ exceptions cannot propagate through M3 stack frames.  The C
bridge catches all exceptions and converts them to error indicators:

\begin{lstlisting}[language=C]
toy_shape_t toy_circle_new(double cx, double cy, double r) {
    try {
        return new Circle(cx, cy, r);
    } catch (const std::exception &e) {
        set_error(e.what());
        return NULL;
    } catch (...) {
        set_error("unknown C++ exception");
        return NULL;
    }
}
\end{lstlisting}

The M3 side checks for \texttt{NULL} and raises an M3 exception:

\begin{lstlisting}
PROCEDURE WrapHandle(h: Handle): Shape RAISES {Error} =
  BEGIN
    IF h = NIL THEN RAISE Error(GetLastError()); END;
    ...
  END WrapHandle;
\end{lstlisting}

The error message is stored in a thread-local buffer on the C side
and retrieved by \texttt{toy\_last\_error()}.

%% -------------------------------------------------------------------
\section{Linking}

The C++ wrapper is compiled as a static library:

\begin{lstlisting}
c++ -std=c++17 -O2 -fPIC -c shapes_c.cpp
ar rcs libtoyshapes.a shapes_c.o
\end{lstlisting}

The M3 \texttt{m3makefile} links it with \texttt{import\_lib} and
adds the C++ runtime via \texttt{SYSTEM\_LIBS}:

\begin{lstlisting}
import_lib("toyshapes", path() & "/../../toylib")
SYSTEM_LIBS{"CXX"} = ["-lc++"]
SYSTEM_LIBORDER += ["CXX"]
import_sys_lib("CXX")
\end{lstlisting}

On macOS, \texttt{-lc++} links the system \texttt{libc++.dylib}.
On Linux, this would be \texttt{-lstdc++}.

%% -------------------------------------------------------------------
\section{Generalizing to Real Libraries}

The patterns demonstrated here apply directly to binding any
C++ library (Open3D, Eigen, CGAL, etc.):

\begin{enumerate}
\item Write a \textbf{C bridge} (\texttt{extern "C"} in a
  \texttt{.cpp} file) that exposes the needed functionality
  through opaque handles, flat arrays, and function-pointer
  callbacks.

\item Write an \textbf{M3 raw interface} (\texttt{UNSAFE INTERFACE})
  that matches the C bridge exactly, with \texttt{EXTERNAL}
  declarations.

\item Write an \textbf{M3 idiomatic layer} that wraps handles in
  traced objects with GC cleaners, converts strings, and provides
  closure-based callbacks.

\item Link the C++ static library and \texttt{-lc++} into the
  M3 program.
\end{enumerate}

The main cost is writing and maintaining the C bridge.  For a large
library like Open3D, this is significant work, but it only needs to
cover the functions actually used.

\subsection{Open questions}

\begin{itemize}
\item \textbf{Shared libraries}: The prototype uses static linking.
  For large C++ libraries, dynamic linking
  (\texttt{libtoyshapes.dylib}) avoids code duplication but requires
  the library to be findable at runtime.

\item \textbf{Thread safety}: The error buffer uses
  \texttt{thread\_local}.  If M3 and C++ both use threads, the
  interaction between the M3 GC (which may stop threads) and C++
  thread primitives needs care.

\item \textbf{Large data}: For million-element arrays (vertices,
  etc.), the copy-in/copy-out pattern works but doubles memory.
  A zero-copy approach could share a buffer allocated by one side
  and accessed by both, but requires careful lifetime management.

\item \textbf{Subclassing from M3}: The current design does not
  allow M3 code to create new Shape subclasses that participate in
  C++ virtual dispatch.  This would require a C++ proxy class
  that delegates virtual calls back to M3 methods---doable but
  complex.
\end{itemize}

\end{document}
